\section{%
Introduction
}

One of the most elegant formulas in combinatorics is the Frame--Robinson--Thrall hook formula~\cite[Theorem~1]{FRT} for the number of standard tableaux.
Given a partition $\lambda$ of $n$, 
a standard tableaux of shape $\lambda$ is a filling of the cells of 
the Young diagram $D(\lambda)$ of $\lambda$ with numbers $1, 2, \dots, n$ 
such that each number appears once and the entries of each row and each column are increasing.
The Frame--Robinson--Thrall hook formula asserts that the number $f^\lambda$ of 
standard tableaux of shape $\lambda$ is given by
\begin{equation}
\label{eq:FRT-hook}
f^\lambda
 =
\frac{ n! }
 { \prod_{v \in D(\lambda)} h_{D(\lambda)}(v) },
\end{equation}
where $h_{D(\lambda)}(v)$ denotes the hook length of the cell $v$ in $D(\lambda)$.
Similar formulas hold for the number of shifted standard tableaux 
(\cite[5.1.4, Exercise~21]{Knu}, see also~\cite[\S~40]{Schur} and~\cite[Theorem~1]{T}) 
and the number of increasing labeling of rooted trees (\cite[5.1.4, Exercise~20]{Knu}).
These tableaux and labelings can be regarded as linear extensions of certain posets.

Stanley~\cite{Stan1} introduced the notion of $P$-partitions for a poset $P$, 
and found a relationship between the univariate generating function 
and the number of linear extensions of $P$.
Given a poset $P$, a $P$-partition is an order-reversing map $\sigma$ from $P$ to $\Nat$, 
the set of nonnegative integers.
We denote by $\AP(P)$ the set of all $P$-partitions.
For a $P$-partition $\sigma$, we write $|\sigma| = \sum_{v \in P} \sigma(x)$.
Then Stanley~\cite[Corollaries~5.3 and~5.4]{Stan1} proved that, for a poset $P$ with $n$ elements, 
there exists a polynomial $W_P(q)$ satisfying
\begin{equation}
\label{eq:P-partition}
\sum_{\sigma \in \AP(P)} q^{|\sigma|}
 =
\frac{ W_P(q) }
 { \prod_{i=1}^n (1 - q^i) },
\end{equation}
and that $W_P(1)$ is equal to the number of linear extensions of $P$.
Also in~\cite[Proposition~18.3]{Stan2} he proved that, 
if $P$ is the Young diagram $D(\lambda)$ of a partition $\lambda$, viewed as a poset, 
the generating function of $D(\lambda)$-partitions (also called reverse plane partitions of shape $\lambda$) 
is given by
\begin{equation}
\label{eq:Stanley-hook}
\sum_{\sigma \in \AP(D(\lambda))} q^{|\sigma|}
 =
\frac{ 1 }
 { \prod_{v \in D(\lambda)} (1 - q^{h_{D(\lambda)}(v)}) }.
\end{equation}
Combining~\eqref{eq:P-partition} and~\eqref{eq:Stanley-hook} and taking the limit $q \to 1$, 
we obtain the Frame--Robinson--Thrall hook formula~\eqref{eq:FRT-hook}.
Gansner~\cite[Theorem~5.1]{G} gave a multivariate generalization of~\eqref{eq:Stanley-hook}.

Proctor~\cite{Proc1, Proc2} introduced a wide class of posets, 
called $d$-complete posets, enjoying ``hook-length property'', 
as a generalization of Young diagrams, shifted Young diagrams and rooted trees.
$d$-Complete posets are defined by certain local structural conditions 
(see Section~2 for a precise definition).
Peterson and Proctor obtained the following theorem, which is a far-reaching generalization 
of the hook formulas~\eqref{eq:FRT-hook} and~\eqref{eq:Stanley-hook}.

\begin{theo}[Peterson--Proctor, see~\cite{Proc3}]
\label{thm:PP}
%\textup{(Peterson--Proctor, see~\cite{Proc3})}
Let $P$ be a $d$-complete poset.
The multivari\-ate generating function of $P$-partitions is given by
\begin{equation}
\label{eq:PP-hook}
\sum_{\sigma \in \AP(P)} \vectz^{\sigma}
 =
\frac{ 1 }
 { \prod_{v \in P} (1 - \vectz[H_P(v)]) }.
\end{equation}
(Refer to Section~2 for undefined notations.)
\end{theo}

However the original proof of this theorem is not yet published, 
though an outline of their proof is given in~\cite{Proc3}.
Different proofs are sketched by Ishikawa--Tagawa~\cite{IT1, IT2} and Nakada~\cite{N3, N2}.
Our skew generalization (Theorem~\ref{thm:main} below) provides an alternate proof of Theorem~\ref{thm:PP}.
In the univariate case, a full proof is given by Kim--Yoo~\cite{KY}.

Another direction of generalizing the Frame--Robinson--Thrall hook formula~\eqref{eq:FRT-hook} 
is to consider skew shapes.
For partitions $\lambda \supset \mu$, 
a standard tableau of skew shape $\lambda/\mu$ is a filling of the cells of the skew Young diagram 
$D(\lambda/\mu) = D(\lambda) \setminus D(\mu)$ satisfying the same conditions as standard tableaux 
of straight shape.
However one cannot expect a nice product formula for the number $f^{\lambda/\mu}$ of 
standard tableaux of skew shape $\lambda/\mu$ in general.
Naruse~\cite{Naruse} presented and sketched a proof of a subtraction-free formula for $f^{\lambda/\mu}$:
\begin{equation}
\label{eq:Naruse-hook}
f^{\lambda/\mu}
 =
n!
\sum_{D \in \EE_{D(\lambda)}(D(\mu))} 
 \frac{ 1 }
 { \prod_{v \in D(\lambda) \setminus D} h_\lambda(v) },
\end{equation}
where $n = |\lambda/\mu|$, 
and $D$ runs over all excited diagrams of $D(\mu)$ in $D(\lambda)$.
Morales--Pak--Panova~\cite{MPP} gave a $q$-analogue of Naruse's skew hook formula 
for the univariate generating functions for $P$-partitions on $P = D(\lambda/\mu)$.

The main result of this paper is the following skew generalization of Peterson--Proctor's 
hook formula (Theorem~\ref{thm:PP}).
Recall that a subset $F$ of a poset $P$ is called an order filter of $P$ if 
$x<y$ in $P$ and $x \in F$ imply $y \in F$.
In particular the empty set $F = \emptyset$ is an order filter of $P$.
